Preventing the spread of nuclear material for weapons is the goal of nuclear safeguards.
%Nuclear safeguards is an important part of the nuclear industry, as preventing the spread of nuclear material for weapons is critical.
%Although there are benefits to \glspl{MSR}, there are many open technical issues that still must be resolved. 
%Within nuclear safeguards there are two general categories: \gls{PR} and \gls{PP} \cite{protection2011evaluation}.
Safeguards, referred to as \gls{PR} here, is defined as the nuclear reactor characteristics that impede "...the diversion or undeclared production of nuclear material and the misuse of technology by the Host State seeking to acquire nuclear weapons or other nuclear explosive devices." \cite{protection2011evaluation}.
This definition agrees with the definition made at the \gls{IAEA} international workshop in 2002 \cite{IAEA2002Proliferation}. 

 There are six measures for \gls{PR}: technical difficulty, cost, time, fissile material type, detection probability, and detection resource efficiency.
 In this work, the measures analyzed include time and detection probability.
 The time measure is analyzed by conducting diversion scenarios over various time frames to determine the effects on diversion fingerprints.
 A very low \gls{PR} time is 0-3 months, while a very high \gls{PR} time is $>$30 years \cite{protection2011evaluation}.
 The detection probability measure goes hand-in-hand with time, as a protracted diversion scenario is likely to have a reduced detection probability.
 A very low \gls{PR} detection probability is 0-5 percent, while a very high \gls{PR} detection probability is 95-100 percent \cite{protection2011evaluation}.

The \gls{DOE} released an updated nuclear \gls{MCA} Order in February of 2023 that cancels \gls{DOE} Order 474.2 Chg 4, which was last updated in 2016 \cite{holmer2024doe}.
Based on the more recent Order, \gls{SNM} is classified as shown in Table \ref{tab:SNM-class}.

\renewcommand{\arraystretch}{1.25}
\begin{table}[H]
    \centering
    \begin{threeparttable}
    \caption{The \Gls{SNM} classification table is recreated from \cite{holmer2024doe}.}
    \label{tab:SNM-class}
    \begin{tabularx}{\textwidth}{X X X X}
        \hline
        Material Type & Accountable Quantity & Weight Field Used for Element & Weight Field Used for Isotope \\
        \hline
        $^{235}$U & 1 g & total uranium & $^{235}$U\\
        $^{233}$U & 1 g & total uranium & $^{233}$U\\
        $^{242}$Pu & 1 g & total plutonium & $^{242}$Pu\\
        $^{239-241}$Pu & 1 g & total plutonium & $^{239}$Pu + $^{241}$Pu\\
        $^{238}$Pu & 0.1 g & total plutonium & $^{238}$Pu\\
        \hline
        $^{241}$Am\tnote{**} & 1 g & total americium & $^{241}$Am\\
        $^{243}$Am\tnote{**} & 1 g & total americium & $^{243}$Am\\
        $^{237}$Np\tnote{**} & 1 g & total neptunium & -\\
        \hline
    \end{tabularx}
    \begin{tablenotes}
      \item[**] Americium and neptunium are treated as \gls{SNM} if they are separated.
    \end{tablenotes}
    \end{threeparttable}
\end{table}
\renewcommand{\arraystretch}{1}

Table \ref{tab:SNM-class} is used along with the graded safeguards table from the updated Order to properly classify the attractiveness level and categorization of \gls{SNM}.
These categories include weapons, pure products, high-grade materials, low-grade materials, and all other materials, ranging from attractiveness level "A" to "E".
For this work, only materials of attractiveness level "C" and below are relevant.
High-grade materials, with attractiveness level "C", include fuel elements and assemblies.
Low-grade materials, with attractiveness level "D", include solutions with target concentrations of 1-25 g/L.
All other materials, with attractiveness level "E", are made up of highly irradiated forms, solutions with concentrations $<$1 g/L, or uranium containing $<$20\% $^{235}$U or $<$10\% $^{233}$U.

The quantities provided in Table \ref{tab:SNM-class} show the minimum reportable, or accountable, quantities for each material type.
These quantities make up the lower limit of Category IV materials for "B" through "E" attractiveness level materials.
The \gls{SNM} Category depends on the quantity of material as well as the attractiveness level of that material.
Higher Category material has more stringent requirements due to their greater strategic importance, with Category I \gls{SNM} having the most requirements.
For example, if material is classified as attractiveness level "C", then any quantity of plutonium or $^{233}$U $\ge$ 6 kg is considered Category I \gls{SNM}.
Alternatively, if the same material contains $^{235}$U, separated $^{237}$Np, or $^{241}$Am and $^{243}$Am, then an amount $\ge$ 20 kg is considered Category I \gls{SNM}.

If the material in the reactor is classified as attractiveness level "D", then any quantity $\ge$ 16 kg is considered Category II, as Category I quantities do not exist for this attractiveness level.
Material of attractiveness level "D" do not have a listed Category I quantity for material containing $^{235}$U, separated $^{237}$Np, or $^{241}$Am and $^{243}$Am.

Material of attractiveness level "E" does not have Category I, II, or III quantities.
Instead, reportable quantities are considered Category IV.


Significant quantities (SQs) of material are also of interest, of which there are two categories: direct use and indirect use.
Direct use nuclear material consists of $^{233}$U, \gls{HEU} ($^{235}$U enriched over 20\%), and any amount of $^{239}$Pu with less than 80\% $^{238}$Pu.
Indirect use nuclear material consists of thorium and \gls{LEU} ($^{235}$U enriched at less than 20\%). 
SQs of the different material categories are given in Table \ref{tab:sqs} \cite{protection2011evaluation, international2022iaea, bathke2012attractive}.

\begin{table}[]
    \centering
    \begin{tabular}{ccr}
        \hline
        Use & Material & Significant Quantity $[kg]$ \\
        \hline
        Direct & Plutonium & 8\\
        Direct & \gls{HEU} & 25\\
        Direct & $^{233}$U & 8\\
        Indirect & $^{233}$U & 75\\
        Indirect & Natural Uranium & 10,000\\
        Indirect & Depleted Uranium & 20,000\\
        \hline
    \end{tabular}
    \caption{Significant quantities of various materials are given along with their use category.}
    \label{tab:sqs}
\end{table}

%\begin{itemize}
%\item Direct use
%\begin{itemize}
%    \setlength\itemsep{-0.25em}
%    \item 8kg of plutonium
%    \item 25kg of \gls{HEU}
%    \item 8kg of $^{233}$U
%\end{itemize}

%\item Indirect use
%\begin{itemize}
%    \setlength\itemsep{-0.25em}
%    \item 75kg of $^{233}$U
%    \item 10t of natural uranium
%    \item 20t of depleted uranium
%\end{itemize}
%\end{itemize}

This information is important for the current work, as it provides values for quantities of diverted material to analyze.
The exact classification, materials, and quantities to use will vary based on the specifics of the fuel salt and diversion scenario.
However, this information provides a set of data which can be directly referenced to determine the diversion rates.

For example, a fuel salt categorized as direct use with a plutonium diversion would have 8kg of plutonium as a significant quantity.
If that fuel salt also had an attractiveness level of "D", then 16kg of plutonium would be Category II.
Then, the diversion time could vary over a range of months to years, based on various \gls{PR} times.
A method of detection for these diversion scenarios is via change in reactivity due to the change in delayed neutron precursors and fissile concentration.
The detection probability is then estimated based on this measurable reactivity change.